\documentclass[aspectratio=169,UTF8,zihao=-4]{ctexbeamer}

\usetheme{Madrid}
\usecolortheme{default}
\setbeamertemplate{navigation symbols}{}
\setbeamertemplate{headline}{}
\setbeamertemplate{frametitle}{}
\setbeamertemplate{footline}[frame number]

\usepackage{amsmath}
\usepackage{booktabs}
\usepackage{array}

\title{范文1第5次提交\\调查过程与样本}
\author{王鲲淏-22307100011}
\date{Social Research}

\newcommand{\bigidea}[1]{
  \begin{center}
    \vspace{0.8em}
    {\LARGE \bfseries #1}
    \vspace{0.6em}
  \end{center}
}

\newcommand{\sectionpageframe}[1]{
  \begin{frame}[plain]
    \begin{center}
      \vfill
      {\Huge \bfseries #1}
      \vfill
    \end{center}
  \end{frame}
}

\begin{document}

\begin{frame}[plain]
  \titlepage
\end{frame}

\sectionpageframe{调查过程}

\begin{frame}
  \bigidea{论文的调查过程}
  \begin{enumerate}
    \item \Large 研究采用 vignette-based survey，每位被试先收到一条经过操纵的 SMS 电影促销广告。
    \item \Large 研究者构造了 16 个情境：$2 \times 2 \times 2 \times 2$，分别操纵折扣、来源、远近和时间。
    \item \Large 次日通过 e-mail 发放问卷，并电话确认被试收到短信、按“刚收到时的反应”填写。
    \item \Large 最后剔除无效和缺失问卷，保留有效样本做后续分析。
  \end{enumerate}
\end{frame}

\begin{frame}
  \bigidea{调查过程中的关键细节}
  \begin{itemize}
    \item \Large 先招募 3 名 student confederates，每人约有 60 名同校学生朋友。
    \item \Large 删除重复联系人后共有 160 名受试者，随机分配到 16 个 scenario，每组约 10 人。
    \item \Large 其中 80 人收到“来自朋友”的短信，80 人收到“来自代表商家的陌生人”的短信。
    \item \Large 研究结束后进行 debrief，解释实验中的情境操纵和欺骗设计。
  \end{itemize}
\end{frame}

\sectionpageframe{样本量}

\begin{frame}
  \bigidea{最后的样本量}
  \begin{center}
    {\Large 最终有效样本量：\alert{N = 141}}
  \end{center}
  \vspace{1em}
  \begin{itemize}
    \item \Large 原始联系样本为 160 人，排除无效问卷后回收 141 份有效问卷。
    \item \Large 性别结构：51 名男性（36.1\%），90 名女性（63.8\%）。
    \item \Large 平均年龄：24.5 岁。
    \item \Large 论文把 demographic statistics 作为控制变量进入假设检验。
  \end{itemize}
\end{frame}

\end{document}
